\section{Other exact-label interfaces}
\label{h:dv:section}

Theorem~\ref{h:dv:finite-density} supplies the strong joint \(L^p\)
density used in both applications. We now separate three additional
conclusions that need not be built into that theorem. A direct
difference-kernel argument proves the measurable-test estimate needed
for selection without using compact-test duality. For spherical
restrictions, joint \(L^1\) convergence identifies the limiting density;
a specified pointwise version works for every full-row-rank matrix,
rather than only almost every matrix.

We retain the notation \(\eta_\mu,\lambda_m,\mathcal E_{z,R}\) of
Section~\ref{h:projection} and the exact norm scale \(\mathcal B\)
in \eqref{h:dv:moment-scale}. The coefficient \(H\) controls mass in
all ambient balls, while \(R\) is the radius of a supporting ball;
neither measure is normalized to a posterior.

\subsection{A direct bound for measurable tests}

\begin{proposition}[A direct inequality for measurable tests]
\label{h:dv:measurable-test}
Under the assumptions of Proposition~\ref{h:dv:ambient-gram}, the
exact image measure satisfies
\begin{equation}
 \eta_\mu(E)\le \mathcal B\,\lambda_m(E)^{1-1/p}
 \label{h:dv:set-bound}
\end{equation}
for every Borel set \(E\) of finite \(\lambda_m\)-measure.  In
particular \(\eta_\mu\ll\lambda_m\), and the same inequality holds on
the \(\lambda_m\)-completion after extending \(\eta_\mu\) to that
completion.  For every \(\lambda_m\)-completed-measurable
\(0\le\varphi\le1\),
\begin{equation}
 \int\varphi\,d\eta_\mu
 \le\mathcal B
 \left(\int_{\mathcal E_{z,R}}\varphi\,d\lambda_m\right)^{1-1/p}.
 \label{h:dv:test-bound}
\end{equation}
An infinite integral on the right is allowed.
\end{proposition}

\begin{proof}
Let \(\rho_\delta=(v_m\delta^m)^{-1}
\mathbf1_{B(0,\delta)}\) be the uniform probability density of the
radius-\(\delta\) ball in \(\mathbb R^m\), and set
\(U_\delta(A,y)=\int\rho_\delta(y-As)\,d\mu(s)\).
After translating the common label variable by \(As_1\), the tuple
kernel in its \(p\)th moment becomes
\[
 K_\delta(v_2,\ldots,v_p)
 =\int_{\mathbb R^m}\rho_\delta(b)
       \prod_{i=2}^p\rho_\delta(b-v_i)\,db,
 \qquad v_i=A(s_i-s_1).
\]
It is nonnegative, and Tonelli's theorem gives the exact normalization
\begin{equation}
 \int_{\mathbb R^{m(p-1)}}K_\delta(v_2,\ldots,v_p)
       \,dv_2\cdots dv_p=1.
 \label{h:dv:difference-kernel}
\end{equation}
In particular the kernel depends on \(A\) only through the differences;
no independence of \(As_1\) from those differences is used.

For a full-rank \(V\), the rows of \(AV\) are independent centered
Gaussian vectors with covariance \(V^{\mathsf T}V\).  Their joint
density has supremum
\((2\pi)^{-m(p-1)/2}\det(V^{\mathsf T}V)^{-m/2}\).
Multiplying this density by the unit-mass kernel
\eqref{h:dv:difference-kernel} bounds the expected tuple kernel by
that supremum.  Expand \(\int U_\delta^p\,d\lambda_m\), discard only
the product-null rank exceptions, and apply
\eqref{h:dv:gram-product}.  This gives
\begin{equation}
 \|U_\delta\|_{L^p(\lambda_m)}\le\mathcal B
 \qquad(\delta>0).
 \label{h:dv:difference-moment}
\end{equation}

The measure \(\lambda_m\) is Borel and finite on compact sets in a
Euclidean space.  The Radon extension theorem
\cite[Corollary~2.11]{StanfordRadonMeasures} applies and gives outer
regularity on Borel sets.  If \(O\) is open, then for
each \((A,As)\in O\), all sufficiently small label perturbations
remain in \(O\).  Fatou's lemma followed by H\"older's inequality gives
\[
 \eta_\mu(O)
 \le\liminf_{\delta\downarrow0}\int_O U_\delta\,d\lambda_m
 \le\mathcal B\,\lambda_m(O)^{1-1/p}.
\]
For a Borel \(E\) with finite \(\lambda_m(E)\), take open supersets
whose \(\lambda_m\)-measures decrease to \(\lambda_m(E)\).  The last
inequality proves \eqref{h:dv:set-bound}.  Applied to null sets, it
gives \(\eta_\mu\ll\lambda_m\).  Extending \(\eta_\mu\) by zero on
\(\lambda_m\)-null sets gives the asserted version on the
\(\lambda_m\)-completion.

The measure \(\eta_\mu\) is supported on the Borel set
\(\mathcal E_{z,R}\).  If
\(\int_{\mathcal E_{z,R}}\varphi\,d\lambda_m<\infty\), apply
\eqref{h:dv:set-bound} to
\(\{\varphi>t\}\cap\mathcal E_{z,R}\) for \(t>0\).  Layer cake and
concavity of \(x^{1-1/p}\) yield
\begin{align*}
 \int\varphi\,d\eta_\mu
 &\le\mathcal B\int_0^1
   \lambda_m(\{\varphi>t\}\cap\mathcal E_{z,R})^{1-1/p}\,dt\\
 &\le\mathcal B
   \left(\int_0^1
   \lambda_m(\{\varphi>t\}\cap\mathcal E_{z,R})\,dt\right)^{1-1/p},
\end{align*}
which is \eqref{h:dv:test-bound}.  The case of an infinite right side
is immediate.  This proof supplies the measurable-test estimate
directly; it does not infer a strong \(L^p\) bound from the set
inequality alone.
\end{proof}

\subsection{Spherical absolute continuity and compatible fibers}

The remaining constructions use an additional hypothesis:
\(\mu\ll\sigma_d\) and \(m<d\). It provides absolute continuity for
each full-row-rank matrix separately. This hypothesis is not needed
by the general finite-measure theorem or the preceding direct-test
bound.

\begin{lemma}[Absolute continuity of spherical projections]
\label{h:dv:gaussian-ac}
Let \(d\ge2\) and \(1\le m<d\).  If \(A:\mathbb R^d\to\mathbb R^m\)
has full row rank, then \(A_\#\sigma_d\) is absolutely continuous with
respect to Lebesgue measure.  The same holds for \(A_\#\mu\) whenever
\(\mu\) is a finite Borel measure with \(\mu\ll\sigma_d\).
\end{lemma}

\begin{proof}
Rotate the row space of \(A\) to the first \(m\) coordinates.  The
restriction of \(A\) to that row space is an invertible map onto
\(\mathbb R^m\), so it suffices to consider the orthogonal projection
onto those coordinates.  Write a uniform sphere point as
\[
 \frac{(Z,W)}{\sqrt{\|Z\|^2+\|W\|^2}},
\]
where \(Z\) and \(W\) are independent standard Gaussian vectors in
\(\mathbb R^m\) and \(\mathbb R^{d-m}\).  Almost surely
\(c=\|W\|>0\).  Conditional on \(c\), the projection is the image of
\(Z\) under
\[
 z\longmapsto\frac{z}{\sqrt{\|z\|^2+c^2}}.
\]
This is a smooth diffeomorphism from \(\mathbb R^m\) onto the open
unit ball; its inverse is
\(u\mapsto cu/\sqrt{1-\|u\|^2}\).  Thus the conditional projection law
has a Lebesgue density, and integrating in \(c\) preserves absolute
continuity.  Applying the invertible row-space map proves the assertion
for \(A_\#\sigma_d\).  Absolute continuity of \(A_\#\mu\) follows from
\(\mu\ll\sigma_d\).
\end{proof}

\begin{proof}[A second proof by hemisphere graphs]
\label{h:dv:hemisphere-ac}
The two open hemispheres of \(S^{d-1}\) are the graphs
\[
 x\longmapsto\bigl(x,\pm\sqrt{1-\|x\|^2}\bigr),
 \qquad \|x\|<1,\quad x\in\mathbb R^{d-1}.
\]
Their surface-area element is
\((1-\|x\|^2)^{-1/2}dx\), and the omitted equator has surface measure
zero.  Projection of \(\sigma_d\) onto the first \(d-1\) coordinates
therefore has a Lebesgue density.  Integrating that density in the last
\(d-1-m\) coordinates proves the same for projection onto the first
\(m\) coordinates.  Rotation, the invertible row-space map, and
\(\mu\ll\sigma_d\) finish the argument as in the first proof.
\end{proof}

\begin{lemma}[A joint density and its fibers]
\label{h:dv:joint-density}
Let \(d\ge2\), \(1\le m<d\), and let \(\mu\) be a finite nonnegative
Borel measure with \(\mu\ll\sigma_d\).  Then \(\eta_\mu\) has a
nonnegative jointly Borel density \(F(A,y)\) with respect to
\(\lambda_m\).  For \(\gamma_m\)-almost every \(A\),
\(F(A,\cdot)\) is a density of \(A_\#\mu\).

If \(\mu\) is supported on \(B(z,R)\), the density can be chosen to
vanish outside the Borel set \(\mathcal E_{z,R}\subseteq\mathcal S_{z,R}\)
defined in \eqref{h:dv:support-sets}.
\end{lemma}

\begin{proof}
For a Borel set \(E\) in the matrix--label space, put
\(E_A=\{y:(A,y)\in E\}\).  If \(\lambda_m(E)=0\), Fubini's theorem
gives \(|E_A|=0\) for \(\gamma_m\)-almost every \(A\).  Such a matrix
has full row rank almost surely, and Lemma~\ref{h:dv:gaussian-ac}
therefore gives \((A_\#\mu)(E_A)=0\).  A second application of Fubini
shows that \(\eta_\mu(E)=0\).  Thus the finite measure \(\eta_\mu\)
is absolutely continuous with respect to the \(\sigma\)-finite measure
\(\lambda_m\).  The ordinary Radon--Nikodym theorem
\cite[Theorem~6.27]{HunterMeasure2011}, applied on the product Borel
\(\sigma\)-algebra, gives a nonnegative jointly Borel density \(F\).

We verify that one common full-measure set of matrices has the asserted
fiber property.  Let \(\mathcal C\) be the countable \(\pi\)-system of
half-open boxes with rational endpoints, the empty set, and
\(\mathbb R^m\) itself.  It generates the Borel sets of
\(\mathbb R^m\).  For
\(C\in\mathcal C\) and every Borel matrix set \(D\), the density
identity applied to \(D\times C\) gives
\[
 \int_D(A_\#\mu)(C)\,d\gamma_m(A)
 =
 \int_D\int_C F(A,y)\,dy\,d\gamma_m(A).
\]
Both integrands are measurable functions of \(A\).  Since this equality
holds for every \(D\), they agree almost everywhere.  Remove the union
of these exceptional sets over the countable class \(\mathcal C\).
For every remaining \(A\), the two finite measures agree on
\(\mathcal C\), and hence on every Borel set by the uniqueness theorem
for measures.  This proves the fiber assertion.

When \(\mu\) is supported on \(B(z,R)\), its image is carried by
the Borel set \(\mathcal E_{z,R}\). Multiplying \(F\) by the indicator
of that set preserves the density identity.
\end{proof}

\subsection{Convergence and an everywhere-full-rank version}

The next result asserts actual convergence in the joint \(L^1\)
space. The final proposition specifies a single Borel version whose
fiber is a density for every full-row-rank matrix, with a separate
Lebesgue-null exceptional set of labels allowed for each matrix.

\Needspace{4\baselineskip}
\begin{proposition}[Joint \(L^1\) convergence]
\label{h:dv:joint-l1}
In addition to the assumptions of
Proposition~\ref{h:dv:ambient-gram}, suppose that \(d\ge2\),
\(1\le m<d\), and \(\mu\ll\sigma_d\).  Let \(F\) be the joint density
from Lemma~\ref{h:dv:joint-density}, and let \(F_\delta\) be the
Gaussian regularization in Lemma~\ref{h:dv:gaussian-moment}.  Then
\begin{equation}
 \|F_\delta-F\|_{L^1(\lambda_m)}\longrightarrow0
 \quad\hbox{as }\delta\downarrow0,
 \qquad
 \|F\|_{L^p(\lambda_m)}\le\mathcal B.
 \label{h:dv:joint-l1-bound}
\end{equation}
\end{proposition}

\begin{proof}
For almost every \(A\), the function \(F_A=F(A,\cdot)\) is a
nonnegative \(L^1\) density of \(A_\#\mu\), of mass
\(\mu(\mathbb R^d)\), and
\(F_\delta(A,\cdot)=F_A*\phi_\delta\).  Tao's continuity theorem for
translations in \(L^1(\mathbb R^m)\)
\cite[Proposition~1.6.13]{TaoMeasure2011} gives
\[
 \|F_A*\phi_\delta-F_A\|_1
 \le\int\phi_\delta(z)\,
          \|F_A(\,\cdot-z)-F_A\|_1\,dz
 \longrightarrow0.
\]
For completeness, fix a small translation radius on which the
translation norm is small.  The integral inside that radius is then
small, and the remaining part is at most
\(2\|F_A\|_1\) times the Gaussian mass outside that radius, which
tends to zero.  This proves the displayed limit.

The fiber difference norm is at most \(2\mu(\mathbb R^d)\).
Because \(F\) and \(F_\delta\) are jointly measurable, these norms
are measurable functions of \(A\).  Dominated convergence in
\(\gamma_m\) therefore proves the joint \(L^1\) limit.  Choose a
sequence \(\delta_i\downarrow0\) for which
\(\sum_i\|F_{\delta_i}-F\|_{L^1(\lambda_m)}<\infty\).
Tonelli's theorem implies \(F_{\delta_i}\to F\) almost everywhere.
Fatou's lemma and \eqref{h:dv:gaussian-moment-bound} then give
\(\int F^p\,d\lambda_m\le\mathcal B^p\).  The density identity from
Lemma~\ref{h:dv:joint-density} now applies to arbitrary bounded
\(\lambda_m\)-completed-measurable functions of the exact
matrix--label pair.
\end{proof}

\begin{proposition}[A pointwise joint density version]
\label{h:dv:pointwise-density}
Assume the same hypotheses as in Proposition~\ref{h:dv:joint-l1}.
For the Gaussian density \(\phi_\delta\), define
\begin{equation}
 w_\delta(r)
 =v_m r^m\frac{r}{\delta^2}\phi_\delta(r e_1),
 \qquad r>0,
 \label{h:dv:radial-mixture}
\end{equation}
where \(e_1\) is a unit coordinate vector.  This is a probability
density on \((0,\infty)\), and \(\phi_\delta\) is the corresponding
mixture of uniform ball densities.  Define
\begin{equation}
 F^\sharp(A,y)=
 \begin{cases}
  \displaystyle\lim_{i\to\infty}F_{1/i}(A,y),
       &\text{if this limit exists and is finite},\\
  0,   &\text{otherwise}.
 \end{cases}
 \label{h:dv:pointwise-version}
\end{equation}
Then \(F^\sharp\) is jointly Borel.  For every full-row-rank \(A\),
\(F^\sharp(A,\cdot)\) is a density of \(A_\#\mu\) up to a Lebesgue
null set, and it is a density of \(\eta_\mu\) with respect to
\(\lambda_m\).  Moreover
\begin{equation}
 \int(F^\sharp)^p\,d\lambda_m\le\mathcal B^p.
 \label{h:dv:pointwise-moment}
\end{equation}
\end{proposition}

\begin{proof}
The negative radial derivative of the Gaussian is
\(-\frac{d}{dr}\phi_\delta(re_1)
=(r/\delta^2)\phi_\delta(re_1)\).  Hence
\[
 \int_{\|y\|}^{\infty}\frac{w_\delta(r)}{v_mr^m}\,dr
 =\phi_\delta(y).
\]
This is precisely the claimed ball mixture.  Integrating it in \(y\)
and using Tonelli proves \(\int_0^\infty w_\delta(r)\,dr=1\).
Also \(w_\delta(r)=\delta^{-1}w_1(r/\delta)\); its probability mass
therefore concentrates at radius zero as \(\delta\downarrow0\).

Fix a full-row-rank \(A\).  Lemma~\ref{h:dv:gaussian-ac} gives a
nonnegative \(L^1\) density \(f_A\) of \(A_\#\mu\).  At a Lebesgue
point \(y\) of \(f_A\), its ball averages
\[
 a_r(y)=\frac{1}{v_mr^m}\int_{B(y,r)}f_A(x)\,dx
\]
tend to \(f_A(y)\) as \(r\downarrow0\), by the Lebesgue differentiation
theorem \cite[Theorem~1.6.19]{TaoMeasure2011}.  The mixture formula gives
\[
 F_\delta(A,y)
 =(f_A*\phi_\delta)(y)
 =\int_0^\infty w_\delta(r)a_r(y)\,dr.
\]
Given an error tolerance, choose \(r_0>0\) such that
\(a_r(y)\) is within that tolerance of \(f_A(y)\) for \(0<r<r_0\).
For \(r\ge r_0\), the averages are bounded by
\(\mu(\mathbb R^d)/(v_mr_0^m)\).  The part of the mixture on
\([r_0,\infty)\) thus tends to zero after subtracting \(f_A(y)\),
while the part on \((0,r_0)\) has the chosen error bound.  It follows
that \(F_\delta(A,y)\to f_A(y)\) at every Lebesgue point.

Each \(F_{1/i}\) is jointly Borel.  The set where its sequence has a
finite limit is Borel, so \eqref{h:dv:pointwise-version} is jointly
Borel as well.  The preceding fixed-\(A\) argument shows that its
fiber is a density for every full-row-rank \(A\), outside a null set
of labels.  Full row rank holds \(\gamma_m\)-almost surely, and
Fubini therefore identifies \(F^\sharp\) with a joint density of
\(\eta_\mu\).  Finally
\((F^\sharp)^p\le\liminf_i F_{1/i}^p\) everywhere by its definition.
Fatou's lemma and \eqref{h:dv:gaussian-moment-bound} prove
\eqref{h:dv:pointwise-moment}.
\end{proof}

The direct-test argument applies to the arbitrary finite
measure in Proposition~\ref{h:dv:ambient-gram}.  The joint \(L^1\) and
pointwise constructions use the additional spherical absolute
continuity, which holds for every restriction
\(\mu=\mathbf1_{B(z,R)}h\,\sigma_d\) with \(0\le h\le1\).
In each case the resulting test or density identity is for exact
observations.
